\subsection{\texorpdfstring{$G_{5,14}$}{G5,14} の node5 sector は全尺度で排除された}
\label{ug25:notes}

新しい確定結果は定理~\texttt{ug25:no-five} である．
二単位枝を $01,23$ とすると，両枝の node 配分はともに $\{4,1,0\}$ に限られる．
support は引き続き $G_{5,14},G_{5,16},G_{5,17}$ の三型・二十六標識であり，
$G_{5,14}$ 全体が排除されたわけではない．

再利用すべき機構は，第二単位枝による matching の伝播である．
第一単位枝を node5 と仮定し，その位置を他単位枝の端 $2$，
低正位置を core $0$ で $3$，core $1$ で $4$ とすれば，
$q_{02}\le7\le q_{03}$ となる．core $2$ の低正位置は $0$ または $1$ の二択で，
外側 pair の符号と反対符号位置の完全表から
\[
 (q_{02},q_{03},q_{12},q_{13})
 =(7,7,3,3)\quad\text{または}\quad(3,7,7,3)
\]
に落ちる．core $1$ では二つの負低行位置の高分母が $(3,3)$ または $(3,7)$ に
固定されるため，その間の正 pair は $LM=8$ または $20$ である．
この全正因子から残る低 column $B$ について線形式を得る．
三十六場合は全て矛盾し，二つの恒等零線形式も逆 entry
$(6B+2)/6$ の非整数性で除かれる．恒等零の場合に整除条件を捨てないこと．

さらに二本の $\{4,1,0\}$ の符号接着により，
定理~\texttt{cm25:central} は中央四高 core の正 matching を $\{01,23\}$ に固定する．
交差 matching を仮定すると，十六の標識付き低正位置配置は四正規形に落ちる．
三つは分母三と node 零側分母七以上の不整合で除かれ，
相互配置 $(f_0,f_2)=(2,0),f_1=f_3=4$ は元 clutching から
局所 row に $-q$ または零の entry を強制し，pair equation に反する．
新しい中心正 matching の下で二本の $\{4,1,0\}$ を同時に正規化することが次の課題である．
各枝の最大 node 位置には必ず分母三が隣接する．
ただし二つの証人分母が同一辺になる可能性があるので，
「相異なる二本の分母三が存在する」とはまだ結論しない．
外側 pair の正負を先に確定し，二つの小分母が一つの三高 core の
負低行位置へ集まれば，補題~\texttt{ug25:factor-test} をそのまま使える．
三高 core に二高 core 用の $\mathcal P(3),\mathcal P(5)$ を流用しない．

$G_{5,16}$ では第二単位枝がないため，本証明は適用されない．
第一単位枝の node5 正規形だけなら $t,A,B$ が無上界に残る．
小さい局所探索で両端同時 completion が見つからなかったことは証明ではなく，
再開時の研究判断に使う必要も薄いので，その探索ログを主たる研究ノートに残さない．
次の有効な分岐は，固定小分母を含む正三高 core の全族分類か，
第三 core の高 pair・node・対向高 $D$ を加えた消去である．
